% --- Archon physics formalization source begin ---
% archon:physics
% archon:covers PhyXMiniProblems/problem_phyx_mini_0277.lean
% archon:source-report reports/phyx_mini/problem_phyx_mini_0277.source.json

\chapter{Physics problem phyx\_mini\_0277}
\label{ch:PhyXMiniProblems_problem_phyx_mini_0277}

\paragraph{Problem source.}
\#\# Physical scenario

As shown in figure.

\#\# Question

What is the phase constant (positive) for SMH with \$a(t)\$ given in figure if the position function \$x(t)\$ has the form \$x = x\_m \textbackslash{}cos(\textbackslash{}omega t + \textbackslash{}phi)\$ and \$a\_s = 4.0\$ m/s\$\textasciicircum{}2\$?

\#\# Answer choices

- A: A: 1.48 rad

- B: B: 1.65 rad

- C: C: 2.24 rad

- D: D: 1.82 rad

\#\# Figure caption (auxiliary; use the image as primary evidence)

The image is a graph plotting acceleration (a) in meters per second squared \textbackslash{}((\textbackslash{}text\{m/s\}\textasciicircum{}2)\textbackslash{}) on the y-axis against time (t) on the x-axis. The graph displays a smooth, curved line starting at a low point, rising to a peak, and then slightly declining. The y-axis is labeled with \textbackslash{}(a\_s\textbackslash{}) at the peak and \textbackslash{}(-a\_s\textbackslash{}) at the trough, indicating positive and negative acceleration values. The x-axis is marked with \textbackslash{}(t\textbackslash{}). Horizontal grid lines are present, assisting in visualizing the changes in acceleration over time. The curve indicates a progressive increase and subsequent minor decrease in acceleration.

\#\# Dataset metadata

Wave Properties; Physical Model Grounding Reasoning, Numerical Reasoning

\paragraph{Recorded answer/context.}
D: D: 1.82 rad

\paragraph{Figure/image path.}
/root/proposal\_for\_physic/science-mango/phyx\_mini\_run/phyx\_data/test\_image/277.png

\paragraph{Formalization target.}
create a compiling Lean file with sorry bodies at `PhyXMiniProblems/problem\_phyx\_mini\_0277.lean`. The Lean declarations must preserve the physical quantities, dimensions or dimensional roles, figure labels, governing-law hypotheses, and final relation expressed by this problem.
Use Mathlib/Physlib names found through LeanExplore where available. If a physics API is missing, introduce faithful local abstractions rather than scalar placeholder aliases.

\begin{theorem}[Physics formalization target]
\label{thm:physics:phyx_mini_0277:target}
\lean{PhyXMiniProblems.ProblemPhyXMini0277.problem_phyx_mini_0277}
\uses{def:physics:phyx-mini-0277:phyxminiproblems-problemphyxmini0277-accelerationgraphoscillator, def:physics:phyx-mini-0277:phyxminiproblems-problemphyxmini0277-matchesprimaryaccelerationgraph, def:physics:phyx-mini-0277:phyxminiproblems-problemphyxmini0277-hasphysicalharmonicparameters, def:physics:phyx-mini-0277:phyxminiproblems-problemphyxmini0277-satisfiescosineharmonicmotion, def:physics:phyx-mini-0277:phyxminiproblems-problemphyxmini0277-answerchoice, def:physics:phyx-mini-0277:phyxminiproblems-problemphyxmini0277-recordedanswerchoice, def:physics:phyx-mini-0277:phyxminiproblems-problemphyxmini0277-matchesanswertonearesthundredth}
The assigned autoformalize agent should translate this physics problem into Lean declarations in the covered file, with theorem and lemma proof bodies written as `by sorry`.
\end{theorem}
\begin{proof}
This is an autoformalization task, not a proof task. Produce statements that can later be proved by the physics prover without weakening the physical model.
\end{proof}

% --- Archon declaration topology begin ---
\section{Declaration topology}
\begin{definition}[Oscillator Length]
  \label{def:physics:phyx-mini-0277:phyxminiproblems-problemphyxmini0277-oscillatorlength}
  \lean{PhyXMiniProblems.ProblemPhyXMini0277.OscillatorLength}
  A signed one-dimensional displacement or displacement amplitude
\end{definition}
\begin{proof}
  This is the declared model definition of Oscillator Length.
\end{proof}
\begin{definition}[Oscillator Time]
  \label{def:physics:phyx-mini-0277:phyxminiproblems-problemphyxmini0277-oscillatortime}
  \lean{PhyXMiniProblems.ProblemPhyXMini0277.OscillatorTime}
  A physical time coordinate
\end{definition}
\begin{proof}
  This is the declared model definition of Oscillator Time.
\end{proof}
\begin{definition}[Angular Frequency]
  \label{def:physics:phyx-mini-0277:phyxminiproblems-problemphyxmini0277-angularfrequency}
  \lean{PhyXMiniProblems.ProblemPhyXMini0277.AngularFrequency}
  Angular frequency, with radians treated as dimensionless
\end{definition}
\begin{proof}
  This is the declared model definition of Angular Frequency.
\end{proof}
\begin{definition}[Oscillator Acceleration]
  \label{def:physics:phyx-mini-0277:phyxminiproblems-problemphyxmini0277-oscillatoracceleration}
  \lean{PhyXMiniProblems.ProblemPhyXMini0277.OscillatorAcceleration}
  A signed one-dimensional acceleration
\end{definition}
\begin{proof}
  This is the declared model definition of Oscillator Acceleration.
\end{proof}
\begin{definition}[Oscillator Jerk]
  \label{def:physics:phyx-mini-0277:phyxminiproblems-problemphyxmini0277-oscillatorjerk}
  \lean{PhyXMiniProblems.ProblemPhyXMini0277.OscillatorJerk}
  A signed one-dimensional jerk, the physical slope of an acceleration--time graph
\end{definition}
\begin{proof}
  This is the declared model definition of Oscillator Jerk.
\end{proof}
\begin{definition}[meters Value]
  \label{def:physics:phyx-mini-0277:phyxminiproblems-problemphyxmini0277-metersvalue}
  \lean{PhyXMiniProblems.ProblemPhyXMini0277.metersValue}
  \uses{def:physics:phyx-mini-0277:phyxminiproblems-problemphyxmini0277-oscillatorlength}
  Read a physical length in metres
\end{definition}
\begin{proof}
  This is the declared model definition of meters Value.
\end{proof}
\begin{definition}[seconds Value]
  \label{def:physics:phyx-mini-0277:phyxminiproblems-problemphyxmini0277-secondsvalue}
  \lean{PhyXMiniProblems.ProblemPhyXMini0277.secondsValue}
  \uses{def:physics:phyx-mini-0277:phyxminiproblems-problemphyxmini0277-oscillatortime}
  Read a physical time in seconds
\end{definition}
\begin{proof}
  This is the declared model definition of seconds Value.
\end{proof}
\begin{definition}[radians Per Second Value]
  \label{def:physics:phyx-mini-0277:phyxminiproblems-problemphyxmini0277-radianspersecondvalue}
  \lean{PhyXMiniProblems.ProblemPhyXMini0277.radiansPerSecondValue}
  \uses{def:physics:phyx-mini-0277:phyxminiproblems-problemphyxmini0277-angularfrequency}
  Read an angular frequency in radians per second
\end{definition}
\begin{proof}
  This is the declared model definition of radians Per Second Value.
\end{proof}
\begin{definition}[meters Per Second Squared Value]
  \label{def:physics:phyx-mini-0277:phyxminiproblems-problemphyxmini0277-meterspersecondsquaredvalue}
  \lean{PhyXMiniProblems.ProblemPhyXMini0277.metersPerSecondSquaredValue}
  \uses{def:physics:phyx-mini-0277:phyxminiproblems-problemphyxmini0277-oscillatoracceleration}
  Read a signed acceleration in metres per second squared
\end{definition}
\begin{proof}
  This is the declared model definition of meters Per Second Squared Value.
\end{proof}
\begin{definition}[meters Per Second Cubed Value]
  \label{def:physics:phyx-mini-0277:phyxminiproblems-problemphyxmini0277-meterspersecondcubedvalue}
  \lean{PhyXMiniProblems.ProblemPhyXMini0277.metersPerSecondCubedValue}
  \uses{def:physics:phyx-mini-0277:phyxminiproblems-problemphyxmini0277-oscillatorjerk}
  Read a signed jerk in metres per second cubed
\end{definition}
\begin{proof}
  This is the declared model definition of meters Per Second Cubed Value.
\end{proof}
\begin{definition}[Graph Axis]
  \label{def:physics:phyx-mini-0277:phyxminiproblems-problemphyxmini0277-graphaxis}
  \lean{PhyXMiniProblems.ProblemPhyXMini0277.GraphAxis}
  The two coordinate axes visible in the primary graph
\end{definition}
\begin{proof}
  This is the declared model definition of Graph Axis.
\end{proof}
\begin{definition}[Axis Quantity]
  \label{def:physics:phyx-mini-0277:phyxminiproblems-problemphyxmini0277-axisquantity}
  \lean{PhyXMiniProblems.ProblemPhyXMini0277.AxisQuantity}
  The physical quantity printed on each coordinate axis
\end{definition}
\begin{proof}
  This is the declared model definition of Axis Quantity.
\end{proof}
\begin{definition}[Acceleration Time Graph]
  \label{def:physics:phyx-mini-0277:phyxminiproblems-problemphyxmini0277-accelerationtimegraph}
  \lean{PhyXMiniProblems.ProblemPhyXMini0277.AccelerationTimeGraph}
  \uses{def:physics:phyx-mini-0277:phyxminiproblems-problemphyxmini0277-oscillatoracceleration, def:physics:phyx-mini-0277:phyxminiproblems-problemphyxmini0277-graphaxis, def:physics:phyx-mini-0277:phyxminiproblems-problemphyxmini0277-axisquantity}
  Labels, units, scale, and grid data read from the primary figure. The horizontal axis has no printed time unit. The vertical label explicitly uses metres and seconds squared; these units are represented by Physlib's LengthUnit and TimeUnit rather than by an untyped string
\end{definition}
\begin{proof}
  This is the declared model definition of Acceleration Time Graph.
\end{proof}
\begin{definition}[Acceleration Graph Oscillator]
  \label{def:physics:phyx-mini-0277:phyxminiproblems-problemphyxmini0277-accelerationgraphoscillator}
  \lean{PhyXMiniProblems.ProblemPhyXMini0277.AccelerationGraphOscillator}
  \uses{def:physics:phyx-mini-0277:phyxminiproblems-problemphyxmini0277-oscillatorlength, def:physics:phyx-mini-0277:phyxminiproblems-problemphyxmini0277-oscillatortime, def:physics:phyx-mini-0277:phyxminiproblems-problemphyxmini0277-angularfrequency, def:physics:phyx-mini-0277:phyxminiproblems-problemphyxmini0277-oscillatoracceleration, def:physics:phyx-mini-0277:phyxminiproblems-problemphyxmini0277-oscillatorjerk, def:physics:phyx-mini-0277:phyxminiproblems-problemphyxmini0277-accelerationtimegraph}
  The physical quantities referred to by the question and acceleration graph. The phase field stores the unknown requested by the problem, but no value is assigned to it here. Position, acceleration, and jerk are independent dimensionful observables connected to the amplitude, angular frequency, and phase only by the governing-law premise below
\end{definition}
\begin{proof}
  This is the declared model definition of Acceleration Graph Oscillator.
\end{proof}
\begin{definition}[Matches Primary Acceleration Graph]
  \label{def:physics:phyx-mini-0277:phyxminiproblems-problemphyxmini0277-matchesprimaryaccelerationgraph}
  \lean{PhyXMiniProblems.ProblemPhyXMini0277.MatchesPrimaryAccelerationGraph}
  \uses{def:physics:phyx-mini-0277:phyxminiproblems-problemphyxmini0277-secondsvalue, def:physics:phyx-mini-0277:phyxminiproblems-problemphyxmini0277-meterspersecondsquaredvalue, def:physics:phyx-mini-0277:phyxminiproblems-problemphyxmini0277-meterspersecondcubedvalue, def:physics:phyx-mini-0277:phyxminiproblems-problemphyxmini0277-accelerationgraphoscillator}
  Primary-image evidence and the stated a\_s calibration. The curve meets the marked extrema a\_s and -a\_s, so its acceleration amplitude equals the scale. At the vertical axis it is one of the four grid intervals above zero and is rising. The rising condition is recorded dimensionfully as positive jerk, not as a premise about the phase
\end{definition}
\begin{proof}
  This is the declared model definition of Matches Primary Acceleration Graph.
\end{proof}
\begin{definition}[Has Physical Harmonic Parameters]
  \label{def:physics:phyx-mini-0277:phyxminiproblems-problemphyxmini0277-hasphysicalharmonicparameters}
  \lean{PhyXMiniProblems.ProblemPhyXMini0277.HasPhysicalHarmonicParameters}
  \uses{def:physics:phyx-mini-0277:phyxminiproblems-problemphyxmini0277-metersvalue, def:physics:phyx-mini-0277:phyxminiproblems-problemphyxmini0277-radianspersecondvalue, def:physics:phyx-mini-0277:phyxminiproblems-problemphyxmini0277-meterspersecondsquaredvalue, def:physics:phyx-mini-0277:phyxminiproblems-problemphyxmini0277-accelerationgraphoscillator}
  Nondegeneracy and the standard positive phase representative. Restricting the phase to one turn is a convention, not the requested numerical answer; the graph ordinate and slope must still determine its quadrant and value
\end{definition}
\begin{proof}
  This is the declared model definition of Has Physical Harmonic Parameters.
\end{proof}
\begin{definition}[Satisfies Cosine Harmonic Motion]
  \label{def:physics:phyx-mini-0277:phyxminiproblems-problemphyxmini0277-satisfiescosineharmonicmotion}
  \lean{PhyXMiniProblems.ProblemPhyXMini0277.SatisfiesCosineHarmonicMotion}
  \uses{def:physics:phyx-mini-0277:phyxminiproblems-problemphyxmini0277-oscillatortime, def:physics:phyx-mini-0277:phyxminiproblems-problemphyxmini0277-metersvalue, def:physics:phyx-mini-0277:phyxminiproblems-problemphyxmini0277-secondsvalue, def:physics:phyx-mini-0277:phyxminiproblems-problemphyxmini0277-radianspersecondvalue, def:physics:phyx-mini-0277:phyxminiproblems-problemphyxmini0277-meterspersecondsquaredvalue, def:physics:phyx-mini-0277:phyxminiproblems-problemphyxmini0277-meterspersecondcubedvalue, def:physics:phyx-mini-0277:phyxminiproblems-problemphyxmini0277-accelerationgraphoscillator}
  The governing simple-harmonic-motion laws in the plus-phase convention stated in the question. The acceleration equation is the second physical time derivative of the position equation, and the jerk equation is its next time derivative. The last field relates the acceleration amplitude to x\_m * omega\textasciicircum{}2. Physlib's HarmonicOscillator.AmplitudePhase instead uses scalar coordinates and the convention A cos (omega t - phi).
\end{definition}
\begin{proof}
  This is the declared model definition of Satisfies Cosine Harmonic Motion.
\end{proof}
\begin{definition}[Answer Choice]
  \label{def:physics:phyx-mini-0277:phyxminiproblems-problemphyxmini0277-answerchoice}
  \lean{PhyXMiniProblems.ProblemPhyXMini0277.AnswerChoice}
  Labels of the four phase choices printed with the problem
\end{definition}
\begin{proof}
  This is the declared model definition of Answer Choice.
\end{proof}
\begin{definition}[radians]
  \label{def:physics:phyx-mini-0277:phyxminiproblems-problemphyxmini0277-answerchoice-radians}
  \lean{PhyXMiniProblems.ProblemPhyXMini0277.AnswerChoice.radians}
  \uses{def:physics:phyx-mini-0277:phyxminiproblems-problemphyxmini0277-answerchoice}
  The phase in radians printed beside each answer label
\end{definition}
\begin{proof}
  This is the declared model definition of radians.
\end{proof}
\begin{definition}[recorded Answer Choice]
  \label{def:physics:phyx-mini-0277:phyxminiproblems-problemphyxmini0277-recordedanswerchoice}
  \lean{PhyXMiniProblems.ProblemPhyXMini0277.recordedAnswerChoice}
  \uses{def:physics:phyx-mini-0277:phyxminiproblems-problemphyxmini0277-answerchoice}
  The answer label recorded in the supplied dataset metadata
\end{definition}
\begin{proof}
  This is the declared model definition of recorded Answer Choice.
\end{proof}
\begin{definition}[Matches Answer To Nearest Hundredth]
  \label{def:physics:phyx-mini-0277:phyxminiproblems-problemphyxmini0277-matchesanswertonearesthundredth}
  \lean{PhyXMiniProblems.ProblemPhyXMini0277.MatchesAnswerToNearestHundredth}
  \uses{def:physics:phyx-mini-0277:phyxminiproblems-problemphyxmini0277-answerchoice}
  Agreement with a phase displayed to the nearest hundredth of a radian
\end{definition}
\begin{proof}
  This is the declared model definition of Matches Answer To Nearest Hundredth.
\end{proof}
\begin{lemma}[initial phase cosine eq neg one fourth]
  \label{lem:physics:phyx-mini-0277:phyxminiproblems-problemphyxmini0277-initial-phase-cosine-eq-neg-one-fourth}
  \lean{PhyXMiniProblems.ProblemPhyXMini0277.initial_phase_cosine_eq_neg_one_fourth}
  \uses{def:physics:phyx-mini-0277:phyxminiproblems-problemphyxmini0277-accelerationgraphoscillator, def:physics:phyx-mini-0277:phyxminiproblems-problemphyxmini0277-matchesprimaryaccelerationgraph, def:physics:phyx-mini-0277:phyxminiproblems-problemphyxmini0277-hasphysicalharmonicparameters, def:physics:phyx-mini-0277:phyxminiproblems-problemphyxmini0277-satisfiescosineharmonicmotion}
  The initial acceleration readout and the acceleration-amplitude law give cos phi = -1/4
\end{lemma}
\begin{proof}
  The conclusion follows from the governing relations and figure readouts named in the statement of initial phase cosine eq neg one fourth.
\end{proof}
\begin{lemma}[phase lies in second quadrant]
  \label{lem:physics:phyx-mini-0277:phyxminiproblems-problemphyxmini0277-phase-lies-in-second-quadrant}
  \lean{PhyXMiniProblems.ProblemPhyXMini0277.phase_lies_in_second_quadrant}
  \uses{def:physics:phyx-mini-0277:phyxminiproblems-problemphyxmini0277-accelerationgraphoscillator, def:physics:phyx-mini-0277:phyxminiproblems-problemphyxmini0277-matchesprimaryaccelerationgraph, def:physics:phyx-mini-0277:phyxminiproblems-problemphyxmini0277-hasphysicalharmonicparameters, def:physics:phyx-mini-0277:phyxminiproblems-problemphyxmini0277-satisfiescosineharmonicmotion}
  The rising acceleration graph gives sin phi > 0. Together with the negative cosine and the positive one-turn convention, this puts the phase in quadrant II
\end{lemma}
\begin{proof}
  The conclusion follows from the governing relations and figure readouts named in the statement of phase lies in second quadrant.
\end{proof}
\begin{lemma}[phase constant eq arccos neg one fourth]
  \label{lem:physics:phyx-mini-0277:phyxminiproblems-problemphyxmini0277-phase-constant-eq-arccos-neg-one-fourth}
  \lean{PhyXMiniProblems.ProblemPhyXMini0277.phase_constant_eq_arccos_neg_one_fourth}
  \uses{def:physics:phyx-mini-0277:phyxminiproblems-problemphyxmini0277-accelerationgraphoscillator, def:physics:phyx-mini-0277:phyxminiproblems-problemphyxmini0277-matchesprimaryaccelerationgraph, def:physics:phyx-mini-0277:phyxminiproblems-problemphyxmini0277-hasphysicalharmonicparameters, def:physics:phyx-mini-0277:phyxminiproblems-problemphyxmini0277-satisfiescosineharmonicmotion}
  On the selected quadrant, the exact phase is the principal inverse cosine
\end{lemma}
\begin{proof}
  The conclusion follows from the governing relations and figure readouts named in the statement of phase constant eq arccos neg one fourth.
\end{proof}
\begin{lemma}[arccos neg one fourth matches choice D]
  \label{lem:physics:phyx-mini-0277:phyxminiproblems-problemphyxmini0277-arccos-neg-one-fourth-matches-choice-d}
  \lean{PhyXMiniProblems.ProblemPhyXMini0277.arccos_neg_one_fourth_matches_choice_D}
  \uses{def:physics:phyx-mini-0277:phyxminiproblems-problemphyxmini0277-matchesanswertonearesthundredth}
  The exact inverse-cosine value rounds to 1.82 rad, displayed as D
\end{lemma}
\begin{proof}
  The conclusion follows from the governing relations and figure readouts named in the statement of arccos neg one fourth matches choice D.
\end{proof}
% --- Archon declaration topology end ---
% --- Archon physics formalization source end ---
